\documentclass[a4paper]{article}
% Español (México) con XeLaTeX: fontspec para fuentes Unicode, polyglossia
% para la división silábica y los nombres del documento en la variante mexicana.
\usepackage{fontspec}
\usepackage{polyglossia}
\setdefaultlanguage[variant=mexican]{spanish}
% Los nombres chinos van como «pinyin (original)»: xeCJK da a los caracteres
% Han su propia fuente; sin ella XeLaTeX los omite con un aviso.
\usepackage{xeCJK}
\setCJKmainfont{FandolSong-Regular.otf}
% polyglossia no traduce los nombres de listings.
\AtBeginDocument{\providecommand{\lstlistingname}{}\renewcommand{\lstlistingname}{Listado}%
  \providecommand{\lstlistlistingname}{}\renewcommand{\lstlistlistingname}{Índice de listados}}
\usepackage{amsmath, amssymb}
\usepackage{graphicx}
\usepackage[margin=2.5cm]{geometry}
\usepackage[most]{tcolorbox}
\usepackage{etoolbox}
\usepackage{listings}
\usepackage{hyperref}
\usepackage{booktabs}
\usepackage{float}
\newtcolorbox{knowledgebox}[1]{enhanced, colback=blue!5!white, colframe=blue!75!black, colbacktitle=blue!75!black, coltitle=white, fonttitle=\bfseries, title={#1}, attach boxed title to top left={yshift=-2mm, xshift=2mm}, boxrule=1pt, sharp corners}
\newtcolorbox{importantbox}[1]{enhanced, colback=yellow!10!white, colframe=yellow!80!black, colbacktitle=yellow!80!black, coltitle=black, fonttitle=\bfseries, title={#1}, sharp corners}
\newtcolorbox{warningbox}[1]{enhanced, colback=red!5!white, colframe=red!75!black, colbacktitle=red!75!black, coltitle=white, fonttitle=\bfseries, title={#1}, sharp corners}
\lstset{language=Python, basicstyle=\ttfamily\small, keywordstyle=\color{blue}, stringstyle=\color{red!60!black}, commentstyle=\color{green!60!black}, breaklines=true, frame=single, numbers=left, numberstyle=\tiny\color{gray}, captionpos=b, extendedchars=true}
\newcommand{\notetitle}{DeepSeek Math V2: Autoverificación y entrenamiento de razonamiento}
\newcommand{\noteauthors}{Elaboradas a partir de materiales públicos del curso}
\newcommand{\notedate}{2025}
\newcommand{\videochannel}{Wudaokou Naishi (五道口纳什)}
\newcommand{\videopublishdate}{2025}
\newcommand{\videoduration}{aproximadamente 20 minutos}
\newcommand{\videourl}{}
\newcommand{\videocoverpath}{cover.jpg}
\newcommand{\repourl}{https://github.com/hqhq1025/ai-course-notes}

% Footer with repo info
\usepackage{fancyhdr}
\pagestyle{fancy}
\fancyhf{}
\fancyfoot[C]{\thepage}
\fancyfoot[R]{\footnotesize\color{gray}\href{\repourl}{AI Course Notes}}
\renewcommand{\headrulewidth}{0pt}
\renewcommand{\footrulewidth}{0pt}

\begin{document}
\begin{titlepage}\centering
{\Large Notas del curso\par}\vspace{1.2cm}{\huge\bfseries \notetitle\par}\vspace{0.8cm}{\large \noteauthors\par}\vspace{0.3cm}{\large \notedate\par}\vspace{1.2cm}
\ifdefempty{\videocoverpath}{}{\includegraphics[width=0.82\textwidth,height=0.45\textheight,keepaspectratio]{\videocoverpath}\par}
\vfill
\begin{tcolorbox}[width=0.9\textwidth, colback=black!2!white, colframe=black!60, sharp corners]
\textbf{Autor o canal del video}: \videochannel\par\textbf{Fecha de publicación}: \videopublishdate\par\textbf{Duración del video}: \videoduration\par
\ifdefempty{\videourl}{}{\textbf{Enlace del video}: \href{\videourl}{\nolinkurl{\videourl}}\par}\par
\textbf{Página del proyecto}: \href{\repourl}{\nolinkurl{\repourl}}\par
\end{tcolorbox}
\end{titlepage}
\tableofcontents\newpage

\section{Introducción}
Esta entrega presenta DeepSeek Math V2 y examina cómo entrenar la capacidad de verification (autoverificación) para que quede dentro del modelo.

\begin{importantbox}{Aportación central}
DeepSeek Math V2 entrena al modelo para que tenga capacidad de \textbf{autoverificación} en el razonamiento matemático: después de generar una solución, el modelo puede verificar por sí mismo si la respuesta es correcta y, si detecta un error, retroceder y corregirlo. En esencia, esto equivale a construir un General Reward Model.
\end{importantbox}

\section{Verification vs. razonamiento tradicional}
\subsection{Comparación con Agentic Workflow}
Anteriormente se presentó el verification workflow en modalidad de razonamiento puro (sin entrenamiento del modelo). El avance de DeepSeek Math V2 consiste en \textbf{interiorizar} esta capacidad en el modelo.

\subsection{Resumen de la sección}
Entrenar la capacidad de verification hasta el interior del modelo es el salto de la ``verificación externa a la ``capacidad interna del modelo.

\newpage
\section{Límites del método: verificable no equivale a fácil de entrenar}
La ventaja de la verification está en que el reward es más explícito, pero eso no resuelve automáticamente todas las dificultades del entrenamiento. Incluso si la respuesta se puede verificar, el modelo aún puede aprender muy lento por exploración insuficiente, baja eficiencia de muestreo o pasos intermedios de mala calidad. Por eso, un reward verificable más bien traslada el problema de «si el reward es confiable» a «si la búsqueda y la optimización son suficientemente eficientes».

\begin{importantbox}{El verdadero valor de las tareas verificables}
Ofrecen a la RL un banco de pruebas más limpio: los investigadores pueden concentrar su esfuerzo principal en el muestreo, la exploración, el credit assignment y la estabilidad del entrenamiento, en lugar de lidiar primero, durante un tiempo prolongado, con preferencias humanas ambiguas.
\end{importantbox}

\subsection{Resumen de la sección}
La verification aumenta la confiabilidad de la señal de reward, pero el éxito del entrenamiento sigue dependiendo de si el algoritmo logra convertir realmente esta «supervisión más limpia» en una mejor política.

\newpage
\section{Síntesis y ampliación}
\begin{enumerate}
  \item DeepSeek Math V2 entrena la capacidad de autoverificación del modelo
  \item En esencia es una implementación de un General Reward Model
  \item Es complementario al verification workflow del modo de razonamiento puro
\end{enumerate}
\end{document}

